\documentclass[11pt]{article}
\usepackage[T1]{fontenc}
\usepackage[margin=2.5cm]{geometry}
\usepackage{amsmath,amssymb,amsthm,booktabs}
\usepackage{hyperref}
\title{The relaxed likelihood:\\
existence, escape and consistency of maximum likelihood\\
on the weakly compactified model}
\author{}
\date{10 November 2026}
\newcommand{\E}{\mathbb E}
\newcommand{\cM}{\mathcal M}
\newcommand{\bM}{\overline{\mathcal M}}
\newcommand{\bX}{\overline{X}}
\newcommand{\bP}{\overline P}
\newcommand{\bL}{\overline L}
\newcommand{\doi}[1]{\href{https://doi.org/#1}{doi:#1}}
\newcommand{\status}[1]{\textsf{[#1]}}
\newtheorem{theorem}{Theorem}
\newtheorem{proposition}[theorem]{Proposition}
\newtheorem{lemma}[theorem]{Lemma}
\newtheorem{conjecture}[theorem]{Conjecture}
\theoremstyle{definition}
\newtheorem{definition}[theorem]{Definition}
\newtheorem{example}[theorem]{Example}
\theoremstyle{remark}
\newtheorem*{remark}{Remark}
\begin{document}
\maketitle

\section*{Results and status}
This note records a discussion of 10 November 2026 (Jonas Moss and Claude).
It proposes one framework for the finite-sample outcomes of maximum
likelihood (ML) and for its consistency, organised around the weak closure of
the model on the compactified sample space. Every claim carries a status tag:
\status{proved here} (short proof in the note), \status{sketch} (argument
given, details unchecked), \status{conjecture}, \status{lit, verified}
(checked in the source text by a literature job), \status{lit, unverified}
(recalled or cited secondhand).

\begin{enumerate}
\item The weak closure $\bM$ of the model in the probability measures on the
  compactified sample space is compact. Each boundary point has a mass
  signature (absolutely continuous, singular, escaped).
\item The \emph{relaxed likelihood} $\bL$ is upper semicontinuous on $\bM$ and
  attains $\sup_\theta L$, possibly $+\infty$ (Proposition~\ref{prop:relaxed}).
  Every finite-sample ML outcome is classified by where $\bL$ is maximised;
  the list is exhaustive (Table~\ref{tab:outcomes}).
\item The \emph{relaxation gap} separates genuine escapes from oscillation.
  An MLE is an exact maximiser inside the model; an escape to $X$ is an exact
  maximiser inside family $X$. Discrete models have no gap.
\item Hellinger convergence decides neither existence nor the likelihood
  value; it governs inference near a boundary point.
\item Defective limits can win on finite samples under correct specification
  (Geyer 1990, Example~4.5) but never break consistency under correct
  specification.
\item Consistency fails through separation (sine family) or domination
  (Ferguson). A block version of Wald's conditions, via Kingman's subadditive
  ergodic theorem, is necessary and sufficient for consistency of
  approximate MLEs (Proposition~\ref{prop:kingman}). This is Perlman's
  theorem \cite{perlman} (\status{lit, verified}, Section~\ref{sec:prior}).
\end{enumerate}

Related material: the mass signature and boundary catalogue in
\texttt{\textasciitilde/Files/research/00\_inactive/ml-existence/notes/}
(\texttt{topology-of-likelihood-boundaries.tex},
\texttt{catalogue-of-likelihood-boundaries.tex},
\texttt{censoring-induced-boundary-selection.tex}); the boundary note
\texttt{boundary-points.tex} in this folder; literature reports
\href{evaluations/2026-11-10-defective-limit-lit-report.md}{TASK-171} and
TASK-173 (relaxation and Wald literature, running).

\section{Setting and mass signature}
Let $X$ be a locally compact separable metric space ($\mathbb R$, an
interval, $\mathbb N$) with $\sigma$-finite dominating measure $\nu$, and
$\cM=\{P_\theta:\theta\in\Theta\}$ a dominated model with fixed density
versions $p_\theta$. Let $\bX=X\cup\{\infty\}$ be the one-point
compactification and regard each $P_\theta$ as a law on $\bX$ without mass
at $\infty$. The space of probability measures on $\bX$ is weakly compact and
metrisable (Prokhorov), so the weak closure $\bM$ of $\cM$ is compact.

For $\bP\in\bM$ write $\bP|_X=\bar p\,\nu+P_s$ with $P_s\perp\nu$. The
\emph{mass signature} is
\[
 (m_{ac},m_s,m_\infty)=\Bigl(\int\bar p\,d\nu,\ P_s(X),\ \bP\{\infty\}\Bigr),
 \qquad m_{ac}+m_s+m_\infty=1.
\]
Proper limit: $(1,0,0)$. Collapse: $m_s>0$. Defective: $m_{ac}<1$ with
$m_\infty>0$. One point at infinity suffices for likelihood purposes since
data never sit there; direction matters once censoring makes the tail
observable.

\section{The relaxed likelihood}
\begin{definition}
For a sample $x=(x_1,\dots,x_n)$ and $\bP\in\bM$,
\[
 \bL(\bP;x)=\sup\Bigl\{\limsup_{k\to\infty}L(\theta_k;x):
 P_{\theta_k}\Rightarrow\bP\Bigr\}\in[0,\infty].
\]
\end{definition}

\begin{proposition}[Relaxed existence]\label{prop:relaxed}
$\bL(\cdot;x)$ is upper semicontinuous on $\bM$, and
$\max_{\bM}\bL(\cdot;x)=\sup_\theta L(\theta;x)$; the maximum is attained.
\status{proved here}
\end{proposition}
\begin{proof}
Upper semicontinuity: let $\bP_j\to\bP$ with $\bL(\bP_j)\ge c_j$. Pick
$\theta_j$ with $d(P_{\theta_j},\bP_j)<1/j$ and
$L(\theta_j)\ge\min(c_j,j)-1/j$; then $P_{\theta_j}\Rightarrow\bP$ and
$\bL(\bP)\ge\limsup_j c_j$. Attainment: a maximising sequence $\theta_k$ has a
weakly convergent subsequence with limit $\bP$, so
$\bL(\bP)\ge\sup L$; the reverse inequality is immediate.
\end{proof}

This is the usual relaxation of a functional on a compactification
(calculus of variations). Existence on $\bM$ is free; the content lies in
locating the maximiser and its value.

\begin{definition}[Relaxation gap, exactness]
For $\bP\in\bM$ with density $\bar p$ on $X$ put
\[
 \Delta(\bP;x)=\log\bL(\bP;x)-\sum_{i=1}^n\log\bar p(x_i)\ \ge 0 .
\]
$\bP$ is \emph{exact} for $x$ if $\Delta(\bP;x)=0$.
\end{definition}

\begin{proposition}[MLE and escape]\label{prop:exact}
An MLE exists iff some maximiser of $\bL(\cdot;x)$ is a member $P_\theta$
with $L(\theta;x)=\bL(P_\theta;x)$. \status{proved here}
\end{proposition}
\begin{proof}
If $\hat\theta$ is an MLE then
$\bL(P_{\hat\theta})\ge L(\hat\theta)=\sup L=\max\bL\ge\bL(P_{\hat\theta})$.
The converse is immediate.
\end{proof}

By analogy, \emph{escape to family $X$} means an exact maximiser that is a
member of the named family $X$; this is what a ``use \texttt{mlX}'' message
asserts. The gap can be positive even at members (the sine family with
$\theta=0$ admitted, Example~\ref{ex:sine}), so membership alone is not
enough.

\begin{lemma}[Sufficient conditions for exactness]\label{lem:exact}
(a) If $p_\theta\to\bar p$ uniformly on a neighbourhood of each $x_i$
whenever $P_\theta\Rightarrow\bP$, then $\bP$ is exact for $x$.
(b) If $X$ is countable with counting measure, every $\bP\in\bM$ is exact
for every sample. \status{proved here}
\end{lemma}
\begin{proof}
(a) is immediate. (b) Each $\{x\}$ is open and closed in $\bX$, so
$P_k\Rightarrow\bP$ implies $P_k\{x\}\to\bP\{x\}$.
\end{proof}

So count families have only proper and defective boundary points and no
gaps; oscillation and spikes need continuous data.

\begin{lemma}[Defective limits under normalisation]\label{lem:defective}
Let $\bP$ be exact with $\bar p=c\,g$, $0<c<1$, $\int g=1$. If $g$ is a model
member or an exact point of $\bM$, then $\bL(\bP;x)=c^n\bL(g;x)<\bL(g;x)$
whenever $\bL(g;x)>0$, so $\bP$ is not a maximiser. \status{proved here}
\end{lemma}

\begin{table}[ht]
\centering
\begin{tabular}{@{}llll@{}}
\toprule
maximiser of $\bL$ & signature & value, gap & outcome\\
\midrule
member, exact & n/a & finite, $\Delta=0$ & MLE (\texttt{Ok})\\
proper boundary, exact & $(1,0,0)$ & finite, $\Delta=0$ & escape, ``use \texttt{mlX}''\\
proper boundary, gap & $(1,0,0)$ & finite, $\Delta>0$ & oscillation; no rescue\\
defective boundary & $m_\infty>0$ & finite & escape to a non-model\\
singular boundary & $m_s>0$ & finite & unattained finite supremum\\
singular or spike & any & $+\infty$ & unbounded likelihood\\
several & mixed & tie & multiple maxima, face ties\\
\bottomrule
\end{tabular}
\caption{Finite-sample outcomes. Exhaustive by Proposition~\ref{prop:relaxed};
the library adds ``budget exhausted'' and ``numerical failure'' for
computations that cannot decide the row.}
\label{tab:outcomes}
\end{table}

\section{Hellinger closure}
Hellinger convergence implies weak convergence, so Hellinger limit points
are weak limit points of signature $(1,0,0)$. Hellinger closures are not
compact, so they cannot support Proposition~\ref{prop:relaxed}. They do not
decide the likelihood value either:

\begin{example}[Vanishing spike]\label{ex:spike}
$f(x)=w\,\sigma^{-1}\varphi((x-\nu)/\sigma)+(1-w)\varphi(x-\mu)$ with
$\nu=x_1$, $w,\sigma\to0$, $w/\sigma\to\infty$. Then
$f\to\varphi(\cdot-\mu)$ in Hellinger (the spike has mass $w$), while
$f(x_1)\to\infty$: signature $(1,0,0)$, $\Delta=\infty$. \status{proved here}
\end{example}

Hierarchy: locally uniform convergence near the data to a density of mass
one implies both Hellinger convergence (Scheff\'e) and exactness; neither
converse holds (sine family; Example~\ref{ex:spike}). Hellinger geometry
belongs to inference near the boundary: DQM, LAN, rates, seams and cusps,
boundary tests (Self and Liang \cite{selfliang}).

\section{Representative examples}

\begin{example}[Oscillation: sine family]\label{ex:sine}
$f_\theta(x)=1+\sin(2\pi\theta x)$ on $[-\tfrac12,\tfrac12]$, $\theta>0$; a
density for every real $\theta$ because the sine is odd.
\begin{enumerate}
\item For $n\ge2$ the MLE almost surely does not exist; $\sup L=2^n$ is
  approached only as $\theta\to\infty$. The $x_i$ are a.s. linearly
  independent over $\mathbb Q$, so $\{\theta x\bmod 1\}$ is dense in the
  torus (continuous Kronecker); attainment needs
  $x_2/x_1=(m'+\frac14)/(m+\frac14)$, a null event. For $n=1$ the
  supremum is attained. \status{proved here}
\item $f_\theta\Rightarrow U$ (Riemann--Lebesgue) at both $\theta\to0$ and
  $\theta\to\infty$; $\bL(U)=2^n$, $L_U=1$, $\Delta=n\log2$.
  \status{proved here}
\item Fisher information: $\cos^2u/(1+\sin u)=1-\sin u$ gives
  $I(\theta)=4\pi^2\int x^2(1-\sin 2\pi\theta x)\,dx=\pi^2/3$. DQM expected
  (the root density is Lipschitz in $\theta$). \status{sketch}
\item Consistent estimators exist for $\theta_0>0$: $\|F_\theta-F_U\|_\infty\le
  1/(\pi\theta)$, the map $\theta\mapsto F_\theta$ is injective and
  continuous, and both ends tend to $F_U\ne F_{\theta_0}$, so minimum
  Kolmogorov distance is consistent. A one-step correction should then be
  efficient. The failure is specific to ML. If $\theta=0$ is admitted, $U$ is
  a member with $\Delta>0$, and no estimator is parameter-consistent at
  $\theta_0=0$. \status{sketch}
\end{enumerate}
The integer-$\theta$ version also works (discrete Kronecker; then even $n=1$
fails). The same mechanism with a countable parameter is published: Le Cam
\cite[\S4, after Ferguson]{lecam} takes densities
$1+\frac12(1-e^{-m})\varphi_{s,m}$ on $(0,1]$, with $\varphi_{s,m}=\pm1$ on a
balanced selection $s$ of the $2^m$ dyadic intervals; density ratios stay in
$[\frac12,2]$ and ML estimates $m=+\infty$ (``essentially another example of
Bahadur''). \status{lit, verified} Exact source of the sine version unknown. The old
\texttt{ml-existence} notes list it as ``oscillation with weak averaging''.
\end{example}

\begin{example}[Unbounded for every sample: three-parameter lognormal]
With a free threshold $\gamma\uparrow x_{(1)}$ the likelihood is unbounded for
every sample (Hill \cite{hill}); only a local maximiser is consistent.
\status{lit, unverified} Three-parameter Weibull and gamma behave the same
way with shape below one; the library already restricts this mechanism
(Fr\'echet $\lambda\ge1$, GEV $\xi\ge-1$; Smith \cite{smith}). Lognormal has
no shape to restrict, so it is the clean representative.
\end{example}

\begin{example}[Defective escape, correct specification: Geyer]\label{ex:geyer}
Geyer \cite[Example 4.5]{geyer} (\status{lit, verified}; journal version
unknown): counting measure on $\{(0,0),(0,-1)\}\cup\{(-k,-l):k,l\ge1\}$,
\[
 P_t(0,0)=\frac{1}{2+e^{-t}},\quad P_t(0,-1)=\frac{e^{-t}}{2+e^{-t}},\quad
 P_t(-k,-l)=\frac{e^{-k\theta_1-lt}}{2+e^{-t}},\quad
 (e^{\theta_1}-1)(e^t-1)=1 .
\]
On the event that all observations equal $(0,0)$, which has positive
probability under every $t$, $L_n(t)=(2+e^{-t})^{-n}\uparrow2^{-n}$ as
$t\to\infty$, where half the mass escapes ($\theta_1\to0$). The normalised
limit $\delta_{(0,0)}$ is not in $\bM$, so Lemma~\ref{lem:defective} does not
apply. The finite-$n$ extension is ours; normalisation checked numerically.
Censored analogues: Birnbaum--Saunders upper face \cite{balzhu} and
defective Gompertz cure models \cite{calsavara}.
A constructed continuous analogue:
$p_\theta=a_\theta\varphi(x)+(1-a_\theta)\varphi(x-\theta-2)$,
$a_\theta=\frac34-\frac14e^{-\theta}$; for samples with all $x_i\le0$ the
supremum is $\prod\frac34\varphi(x_i)$, unattained (ml-existence catalogue).
\end{example}

\begin{example}[Exact but inconsistent: Ferguson]\label{ex:ferguson}
Ferguson \cite{ferguson}; the form below is confirmed by Le Cam's
description \cite{lecam} (\status{lit, verified}), which leaves $\delta$
unspecified (``continuous, decreasing from $1$ to $0$, fast enough''). The
specific $\delta$ is as recalled (\status{lit, unverified}); Le Cam notes
that ``all of Wald's conditions, except one, are satisfied'', matching the
analysis below: $\theta\in[0,1]$,
\[
 f(x\mid\theta)=(1-\theta)\frac{1}{\delta(\theta)}
 \Bigl(1-\frac{|x-\theta|}{\delta(\theta)}\Bigr)^{+}+\frac\theta2\mathbf 1_{[-1,1]}(x),
 \qquad \delta(\theta)=(1-\theta)e^{1-(1-\theta)^{-4}} .
\]
Finite samples: $\theta\mapsto f(x\mid\theta)$ is continuous on $[0,1]$ for
$x<1$, the parameter space is compact and $\theta\mapsto P_\theta$ is a
homeomorphism onto $\cM$, so $\bM=\cM$, every point is exact and the MLE
always exists (\texttt{Ok}). \status{proved here}

Asymptotics: at $\theta=1$ the relaxed density (Section~\ref{sec:consistency})
is $\bar p=\frac12=p_1$, so the point is exact and separated. Domination
fails: for every neighbourhood $U$ of $1$,
$\E_0\sup_{\theta\in U}\log f(X\mid\theta)\gtrsim\int^1(1-x)^{-4}\,dx=\infty$.
\status{sketch}
\end{example}

\begin{proposition}[Ferguson threshold]\label{prop:ferguson-a}
With $\delta_a(\theta)=(1-\theta)\exp(1-(1-\theta)^{-a})$ and $\theta_0>0$:
for $a<1$ Wald's conditions hold and ML is consistent; for $a>1$ ML is
inconsistent, since $1-M_n\asymp n^{-1}$ makes the spike gain
$(1-M_n)^{-a}/n\asymp n^{a-1}\to\infty$. The case $a=1$ is open.
\status{sketch}
\end{proposition}

\section{Consistency}\label{sec:consistency}
Let $P_0=P_{\theta_0}$ be the truth. For $\bP\in\bM$ define the relaxed
density of one observation
\[
 \bar p(x)=\lim_{U\downarrow\bP}\ \sup_{\theta:\,P_\theta\in U}p_\theta(x).
\]
By Jensen, $\E_0\log(\bar p/p_0)\le\log\int\bar p$. Wald's argument
\cite{wald} needs, at every $\bP\ne P_0$:
\begin{itemize}
\item \emph{separation}: $\E_0\log(\bar p/p_0)<0$, implied by
  $\int\bar p\le1$ with $\bar p\ne p_0$;
\item \emph{domination}: $\E_0\sup_{U}\log p_\theta(X)<\infty$ for some
  neighbourhood $U$ (then monotone convergence passes to $\bar p$).
\end{itemize}
By Fatou a pointwise limit along one path integrates to at most one, so
separation fails only when $\bar p$ mixes paths (a gap). Under correct
specification defective limits never break consistency:
$\mathrm{KL}(P_0\,\|\,cg)\ge-\log c>0$.

\begin{center}
\begin{tabular}{@{}lll@{}}
\toprule
condition & fails in & mechanism\\
\midrule
separation & sine family & $\bar p\equiv2$, jointly achievable (Kronecker)\\
domination & Ferguson; Le Cam's version of Bahadur \cite[\S4]{lecam} & spikes, mass-integrable, log-nonintegrable\\
\bottomrule
\end{tabular}
\end{center}

Le Cam's version of Bahadur's example (\status{lit, verified} in TASK-171)
selects a defective density limit asymptotically, with the lost mass
concentrating at a point (signature $(c,1-c,0)$), not escaping.

\subsection{Wald's conditions are sufficient only}
\begin{enumerate}
\item \emph{Per-observation domination can fail while ML is consistent.}
  Normal location-scale near collapse: $\sup_U\log p_\theta(x)=\infty$, but
  for a pair
  $\sup_{\mu,\sigma}\varphi_\sigma(x_1-\mu)\varphi_\sigma(x_2-\mu)\lesssim
  |x_1-x_2|^{-2}$ and the supremum over small $U$ tends to $0$.
  \status{sketch}
\item \emph{Jensen is crude}: separation can hold with $\int\bar p>1$.
\item \emph{Failed separation alone does not give inconsistency}: $\bar p$
  lets each observation choose its own $\theta$; the sine family is
  inconsistent because the choices are jointly achievable.
\end{enumerate}

\subsection{A block criterion}
For an open $U\subset\bM$ put $\Theta_U=\{\theta:P_\theta\in U\}$ and
\[
 Z_{s,t}(U)=\log\sup_{\theta\in\Theta_U}\prod_{s<i\le t}p_\theta(X_i)
 -\sum_{s<i\le t}\log p_0(X_i).
\]
Then $Z_{0,n+m}\le Z_{0,n}+Z_{n,n+m}$, and the array is stationary under
shifts of the i.i.d.\ sequence.

\begin{proposition}[Wald with blocks; proposed iff]\label{prop:kingman}
Assume separability (suprema over a countable dense set) and
$\E_0 Z_{0,m}(U)^+<\infty$ for some $m$. Kingman's subadditive ergodic
theorem \cite{kingman} gives, almost surely,
\[
 \gamma(U)=\lim_{n\to\infty}\frac1n Z_{0,n}(U)
 =\inf_{m\ge1}\frac1m\,\E_0 Z_{0,m}(U)\in[-\infty,\infty).
\]
Put $\gamma(\bP)=\inf_{U\ni\bP}\gamma(U)$ and $\gamma^*=\max_{\bM}\gamma$.
Then every approximate MLE (slack $o(n)$ in log-likelihood, a.s.) converges
weakly to $P_0$ iff $\gamma(\bP)<\gamma^*$ for every $\bP\ne P_0$. When
$\gamma(P_0)=0$ this reads: every $\bP\ne P_0$ has a neighbourhood $U$ and a
block size $m$ with
\[
 \frac1m\,\E_0\log\sup_{\theta\in\Theta_U}L_m(\theta)<\E_0\log p_0 .
\]
\status{sketch; essentially Perlman \cite[Thm.~2.2(ii), Thm.~4.1]{perlman}, see Section~\ref{sec:prior}}
\end{proposition}
\begin{proof}[Sketch]
Sufficiency: cover $\bM$ minus a neighbourhood of $P_0$ by finitely many
$U_j$ with $\gamma(U_j)<\gamma^*$; eventually every $\Theta_{U_j}$ loses to the
global supremum by a linear margin. Necessity: if $\gamma(\bP)=\gamma^*$ for
some $\bP\ne P_0$, then $\sup_{\Theta_U}\ell_n\ge\sup\ell_n-o(n)$ for every
$U\ni\bP$, and a diagonal choice over shrinking $U_j$ gives an approximate
MLE converging to $\bP$.
\end{proof}

Points to inspect: (i) whether $\frac1n\sup\ell_n-\frac1n\ell_n(\theta_0)\to
\gamma^*$ holds as claimed (finite cover plus upper semicontinuity of
$\gamma$); (ii) the integrability hypothesis, and whether $\gamma=+\infty$
follows when it fails for all $m$; (iii) exact MLE versus slack $o(n)$ at
ties $\gamma(\bP)=\gamma^*$; (iv) weak versus parameter consistency when a
boundary point coincides with a member; (v) prior art, especially Perlman
\cite{perlman}.

Checks: normal collapse $\gamma=-\infty$ with $m=2$; Cauchy location-scale
collapse: $p(x_1)p(x_2)\to(\pi d)^{-2}$ with $d=|x_1-x_2|$, while three
observations give $\sim\sigma\to0$, so $m\le3$; sine family $\gamma(U)=\log2
-\E_0\log p_0>0$ near $\theta=\infty$; Ferguson $\gamma=+\infty$ near
$\theta=1$; Ferguson with $a<1$: $m=1$ suffices. \status{sketch}

\begin{conjecture}[Block size and existence]
At a collapse stratum the minimal block size $m^*$ equals the smallest
sample size for which the MLE exists almost surely (normal $2$; Cauchy
location-scale $3$, matching $2k<n$).
\end{conjecture}

\subsection{Prior art: Perlman 1972}\label{sec:prior}
Checked in the source text on 10 November 2026 (\status{lit, verified}).
Perlman \cite{perlman} works with
$u(x,\theta;\theta_0)=\log f(x,\theta)-\log f(x,\theta_0)$ on the complement
$\Omega_r$ of a shrinking neighbourhood of $\theta_0$, without compactifying.
\begin{itemize}
\item Definition~1: $u$ is \emph{dominated (by $0$)} on $\Gamma$ if for some
  $k$, $\sup_\Gamma u_k\le s(x_1,\dots,x_k)$ with $\E s<\infty$ ($<0$). This is
  Wald with blocks of size $k$, after Kiefer and Wolfowitz \cite{kw} and
  Berk.
\item Theorem~2.2(ii): if $u$ is dominated and the suprema are measurable,
  then $\limsup_n\sup_\Gamma\frac1n\sum u=\lim_n\frac1n\E\sup_\Gamma\sum_{i\le n}u$
  almost surely. Proof by Berk's reverse martingale over $k$-subsets, not
  Kingman; the content is the same as our $\gamma(U)$.
\item Assumption $\mathcal D$ (dominance for some $k$) is exactly our
  integrability hypothesis, point (ii) above. Under $\mathcal D$, Condition~1
  (dominance by $0$) is necessary and sufficient for the uniform condition
  (1.7), and by Theorem~4.1 (local regularity: local compactness, local
  dominance, upper semicontinuity, measurability) for strong consistency of
  all approximate MLEs.
\item Without $\mathcal D$: Example~2 (sample median, $w=-|x-\theta|$) is
  consistent although not dominated; Perlman handles it by semidominance
  (divide by $b(\theta)$, after Huber \cite{huber}) and, in Section~3, by a
  uniformity condition on $\sigma$-compact sets. So the hole in point (ii)
  is real and known.
\item Perlman warns that the one-point compactification of $\Theta$ is often
  unsuitable (his Example~4:
  $f(x,\theta)\propto1+\frac{a\theta}{1+\theta}\sin(x-\theta)$ on
  $[0,2\pi]$, $0<a<1$; consistent, but Wald, Kiefer--Wolfowitz and Huber
  fail). Our compactification is of the model in the laws on $\bX$, not of
  $\Theta$: there the $\theta\to\infty$ end is the circle of proper members
  $1+a\sin(x-\varphi)$, exact by Lemma~\ref{lem:exact}(a).
\item Strasser \cite{strasser} cites Perlman's criterion as necessary and
  sufficient and improves its measurability side.
\end{itemize}
What remains ours: the finite-sample relaxed likelihood and outcome table,
the mass signature, the gap, and the reading of $\gamma$ as the asymptotic
relaxed likelihood. Prior art for these is the subject of TASK-173.

\section{Rates}
Rates are local and governed by Hellinger geometry. Given consistency, a
regular interior point (DQM, nonsingular information) gives root-$n$
efficiency. A nonregular point with $h^2(\theta_0,\theta_0+u)\asymp|u|^\alpha$
has optimal rate $n^{-1/\alpha}$ (Ibragimov and Has'minskii \cite{ih}); ML
typically attains the rate without efficiency. Near boundary points: seams
(chi-bar-squared), singular-information corners ($n^{1/4}$), cusps. This
is TASK-102.

\section{Two levels}
\begin{center}
\begin{tabular}{@{}lll@{}}
\toprule
level & object & decides\\
\midrule
model & $\bM$, signature, members that coincide with boundary points & whether any consistent estimator exists \cite{lecamschwartz}\\
ML, finite $n$ & $\bL$, $\Delta$, membership & MLE, escape, failure (Table~\ref{tab:outcomes})\\
ML, $n\to\infty$ & $\bar p$, $\gamma(U)$ & consistency (Proposition~\ref{prop:kingman})\\
local & Hellinger, DQM & rates and inference\\
\bottomrule
\end{tabular}
\end{center}

\begin{center}
\begin{tabular}{@{}lll@{}}
\toprule
example & finite sample & asymptotics\\
\midrule
Ferguson & \texttt{Ok}, exact & domination fails\\
sine family & gap, no MLE & separation fails\\
three-parameter lognormal & unbounded & ML undefined; local root consistent\\
Geyer 4.5 & exact defective escape, positive probability & event vanishes\\
vanishing spike & unbounded, Hellinger limit proper & domination fails\\
\bottomrule
\end{tabular}
\end{center}

\section{Open items}
\begin{enumerate}
\item Integrability-free version of Proposition~\ref{prop:kingman}:
  Perlman's semidominance and Section~3 uniformity; Strasser's measurability.
\item Ferguson: fetch the paper, confirm $\delta$, settle $a=1$.
\item Library proof obligations: exactness (Lemma~\ref{lem:exact}) at each
  escape face before ``use \texttt{mlX}''; defective limits dismissed by
  Lemma~\ref{lem:defective}.
\item Prior art for relaxed likelihood and exhaustive classification
  (TASK-173); a journal source for a complete-data defective escape.
\item Conjecture on block sizes.
\end{enumerate}

\begin{thebibliography}{99}
\bibitem{balzhu} N. Balakrishnan and X. Zhu (2014). On the existence and
uniqueness of the maximum likelihood estimates of the parameters of
Birnbaum--Saunders distribution based on Type-I, Type-II and hybrid
censored samples. \emph{Statistics} 48(5), 1013--1032.
\doi{10.1080/02331888.2013.800069}.
\bibitem{calsavara} V. F. Calsavara, A. S. Rodrigues, R. Rocha,
V. F. Tomazella and F. Louzada (2019). Defective regression models for cure
rate modeling with interval-censored data. \emph{Biometrical Journal}
61(4), 841--859. \doi{10.1002/bimj.201800056}.
\bibitem{chen} J. Chen (2017). Consistency of the MLE under mixture models.
\emph{Statistical Science} 32(1), 47--63. Preprint
\href{https://arxiv.org/abs/1607.01251}{arXiv:1607.01251}.
\bibitem{ferguson} T. S. Ferguson (1982). An inconsistent maximum likelihood
estimate. \emph{Journal of the American Statistical Association} 77(380),
831--834. \doi{10.1080/01621459.1982.10477894}.
\bibitem{geyer} C. J. Geyer (1990). \emph{Likelihood and Exponential
Families}. PhD dissertation, University of Washington.
\href{https://www.stat.umn.edu/geyer/thesis/}{stat.umn.edu/geyer/thesis}.
\bibitem{hill} B. M. Hill (1963). The three-parameter lognormal
distribution and Bayesian analysis of a point-source epidemic.
\emph{Journal of the American Statistical Association} 58, 72--84.
(Details unverified.)
\bibitem{ih} I. A. Ibragimov and R. Z. Has'minskii (1981).
\emph{Statistical Estimation: Asymptotic Theory}. Springer.
\bibitem{kingman} J. F. C. Kingman (1973). Subadditive ergodic theory.
\emph{Annals of Probability} 1(6), 883--899. \doi{10.1214/aop/1176996798}.
\bibitem{lecam} L. Le Cam (1990). Maximum likelihood: an introduction.
\emph{International Statistical Review} 58(2), 153--171.
\doi{10.2307/1403464}.
\bibitem{lecamschwartz} L. Le Cam and L. Schwartz (1960). A necessary and
sufficient condition for the existence of consistent estimates.
\emph{Annals of Mathematical Statistics} 31, 140--150. (Details
unverified.)
\bibitem{perlman} M. D. Perlman (1972). On the strong consistency of
approximate maximum likelihood estimators. \emph{Proceedings of the Sixth
Berkeley Symposium on Mathematical Statistics and Probability} 1, 263--281.
\bibitem{strasser} H. Strasser (1981). Consistency of maximum likelihood and
Bayes estimates. \emph{Annals of Statistics} 9(5), 1107--1113.
\doi{10.1214/aos/1176345590}.
\bibitem{kw} J. Kiefer and J. Wolfowitz (1956). Consistency of the maximum
likelihood estimator in the presence of infinitely many incidental
parameters. \emph{Annals of Mathematical Statistics} 27(4), 887--906.
\doi{10.1214/aoms/1177728066}.
\bibitem{huber} P. J. Huber (1967). The behavior of maximum likelihood
estimates under nonstandard conditions. \emph{Proceedings of the Fifth
Berkeley Symposium} 1, 221--233.
\bibitem{selfliang} S. G. Self and K.-Y. Liang (1987). Asymptotic properties
of maximum likelihood estimators and likelihood ratio tests under
nonstandard conditions. \emph{Journal of the American Statistical
Association} 82(398), 605--610. \doi{10.1080/01621459.1987.10478472}.
\bibitem{smith} R. L. Smith (1985). Maximum likelihood estimation in a class
of nonregular cases. \emph{Biometrika} 72(1), 67--90.
\doi{10.1093/biomet/72.1.67}.
\bibitem{wald} A. Wald (1949). Note on the consistency of the maximum
likelihood estimate. \emph{Annals of Mathematical Statistics} 20(4),
595--601. \doi{10.1214/aoms/1177729952}.
\end{thebibliography}
\end{document}
